\section*{Capítulo \ref{ch:g12:adaptive-immunity} --- A imunidade adaptativa e a vacinação}

\begin{solution}{exo:g12:adaptive-immunity:1}
Um \omterm{def:g12:adaptive-immunity:antigen}{antígeno} é qualquer molécula que o sistema adaptativo consegue
reconhecer; um \omterm{def:g12:adaptive-immunity:antigen}{linfócito} é um glóbulo branco que carrega receptores
para um \omterm{def:g12:adaptive-immunity:antigen}{antígeno}. A especificidade dele vem do receptor dele, montado
ao acaso durante a maturação e idêntico em todas as cópias dele.
\end{solution}

\begin{solution}{exo:g12:adaptive-immunity:2}
O \omterm{def:g12:adaptive-immunity:antigen}{antígeno} se liga aos poucos \omterm{def:g12:adaptive-immunity:antigen}{linfócitos} cujo receptor se encaixa;
essas \omterm{def:g10:cells-common-unit:cell}{células} se multiplicam num clone ao longo de dias; o clone se
diferencia em \omterm{def:g10:cells-common-unit:cell}{células} efetoras que combatem; uma minoria vira \omterm{prop:g12:adaptive-immunity:selection}{células
de memória} que persistem.
\end{solution}

\begin{solution}{exo:g12:adaptive-immunity:3}
Uma \omterm{def:g11:gene-expression:protein}{proteína} em forma de Y secretada pelos plasmócitos, uma cópia
solúvel do receptor do \omterm{def:g12:adaptive-immunity:antigen}{linfócito} B; cada braço se liga ao \omterm{def:g12:adaptive-immunity:antigen}{antígeno}.
Numa bactéria ele forma um revestimento que os \omterm{def:g12:innate-immunity:phagocytosis}{fagócitos} agarram, e
amarra bactérias em complexos que são englobados.
\end{solution}

\begin{solution}{exo:g12:adaptive-immunity:4}
Os \omterm{prop:g12:adaptive-immunity:tcells}{linfócitos T auxiliares}, ativados pelas \omterm{prop:g12:innate-immunity:handover}{células dendríticas},
secretam as citocinas que autorizam e amplificam os \omterm{def:g12:adaptive-immunity:antigen}{linfócitos} B e os
outros \omterm{def:g12:adaptive-immunity:antigen}{linfócitos} T. Os \omterm{def:g12:adaptive-immunity:antigen}{linfócitos} T citotóxicos matam as \omterm{def:g10:cells-common-unit:cell}{células} que
exibem fragmentos de um vírus ou de uma \omterm{def:g11:gene-expression:protein}{proteína} de \omterm{def:g11:cancer:cancer}{tumor}.
\end{solution}

\begin{solution}{exo:g12:adaptive-immunity:5}
Uma forma inofensiva do \omterm{def:g12:adaptive-immunity:antigen}{antígeno} do patógeno (atenuado, morto, uma
\omterm{def:g11:gene-expression:protein}{proteína}, uma toxina inativada, ou as instruções dele) com um
adjuvante. Ela provoca uma resposta primária e deixa \omterm{prop:g12:adaptive-immunity:selection}{células de
memória}, de modo que o patógeno real encontre uma resposta secundária.
\end{solution}

\begin{solution}{exo:g12:adaptive-immunity:6}
Primária: uma semana de atraso, um pico baixo (cerca de um décimo do da
secundária) em duas a três semanas, um declínio em um mês.
Secundária: dois ou três dias de atraso, um pico dez vezes mais alto em
dez dias, que persiste por meses.
\end{solution}

\begin{solution}{exo:g12:adaptive-immunity:7}
A memória é específica: as \omterm{prop:g12:adaptive-immunity:selection}{células de memória} deixadas pela primeira
dose de A respondem só a A; B, encontrado pela primeira vez, recebe uma
resposta primária.
\end{solution}

\begin{solution}{exo:g12:adaptive-immunity:8}
Os \omterm{prop:g12:adaptive-immunity:antibodies}{anticorpos} são \omterm{def:g11:gene-expression:protein}{proteínas} dos líquidos e não conseguem entrar nas
\omterm{def:g10:cells-common-unit:cell}{células}. Uma \omterm{def:g10:cells-common-unit:cell}{célula} infectada exibe fragmentos do vírus na superfície
dela, e os \omterm{def:g12:adaptive-immunity:antigen}{linfócitos} T citotóxicos os reconhecem e matam a \omterm{def:g10:cells-common-unit:cell}{célula}.
\end{solution}

\begin{solution}{exo:g12:adaptive-immunity:9}
$1 - 1/4 = 75\%$; $1 - 1/18 \approx 94\%$.
\end{solution}

\begin{solution}{exo:g12:adaptive-immunity:10}
Acima do limiar, cada caso infecta menos de uma outra pessoa, e assim o
vírus não consegue circular e raramente chega a um bebê; os adultos
imunes são a muralha em volta deles.
\end{solution}

\begin{solution}{exo:g12:adaptive-immunity:11}
Cerca de 750, 520 e 180 por microlitro; a \omterm{prop:g12:adaptive-immunity:hiv}{aids} começa por volta do ano
9, quando a contagem cruza 200.
\end{solution}

\begin{solution}{exo:g12:adaptive-immunity:12}
As citocinas dos \omterm{prop:g12:adaptive-immunity:tcells}{linfócitos T auxiliares} são o que autoriza os
\omterm{def:g12:adaptive-immunity:antigen}{linfócitos} B a virar plasmócitos e os \omterm{def:g12:adaptive-immunity:antigen}{linfócitos} T citotóxicos a se
multiplicar; os dois braços dependem do mesmo coordenador, e assim
removê-lo silencia os dois.
\end{solution}

\begin{solution}{exo:g12:adaptive-immunity:13}
Os \omterm{prop:g12:adaptive-immunity:antibodies}{anticorpos} foram selecionados contra o vírus como ele era; o vírus
faz mutar as \omterm{def:g11:gene-expression:protein}{proteínas} de superfície dele mais rápido do que clones
novos conseguem ser selecionados, e assim cada resposta é ultrapassada.
E o vírus se esconde e se multiplica justamente dentro das \omterm{def:g10:cells-common-unit:cell}{células}
auxiliares que coordenariam a resposta.
\end{solution}

\begin{solution}{exo:g12:adaptive-immunity:14}
As \omterm{prop:g12:adaptive-immunity:selection}{células de memória} reconhecem as \omterm{def:g11:gene-expression:protein}{proteínas} de superfície que a
\omterm{def:g12:adaptive-immunity:vaccine}{vacina} apresentou. As da gripe sofrem \omterm{def:g11:mutations:mutation}{mutação} todo ano em formas que a
memória antiga não liga; as \omterm{def:g11:gene-expression:protein}{proteínas} de superfície do sarampo quase
não mudam, e assim a memória de 1970 ainda se encaixa no vírus de hoje.
\end{solution}

\begin{solution}{exo:g12:adaptive-immunity:15}
O que varia: os receptores dos \omterm{def:g12:adaptive-immunity:antigen}{linfócitos}, feitos ao acaso. O que
seleciona: o \omterm{def:g12:adaptive-immunity:antigen}{antígeno}, que deixa se multiplicar apenas os clones que se
encaixam. O que é herdado: o receptor, pelos descendentes do clone,
incluindo as \omterm{prop:g12:adaptive-immunity:selection}{células de memória}. Escala de tempo: dias para uma
resposta, uma vida inteira para a memória, contra gerações para uma
\omterm{def:g12:selection-drift-speciation:population}{população}. A analogia para na transmissão: os clones selecionados
morrem com o corpo e não são passados à geração seguinte.
\end{solution}

\begin{solution}{pb:g12:adaptive-immunity:1}
\textbf{1.} Nada por cerca de uma semana, um pico modesto em duas a
três semanas, um declínio ao longo das semanas seguintes.

\textbf{2.} \omterm{prop:g12:adaptive-immunity:antibodies}{Anticorpos} em dois ou três dias, um pico dez vezes mais
alto, que dura meses: as \omterm{prop:g12:adaptive-immunity:selection}{células de memória} deixadas pela primeira dose
já são um clone grande, e assim a multiplicação parte de milhares de
\omterm{def:g10:cells-common-unit:cell}{células} em vez de algumas.

\textbf{3.} Um vírus vivo se multiplica por pouco tempo dentro das
\omterm{def:g10:cells-common-unit:cell}{células}, e assim as \omterm{def:g10:cells-common-unit:cell}{células} infectadas exibem fragmentos dele e os
\omterm{def:g12:adaptive-immunity:antigen}{linfócitos} T citotóxicos são selecionados além dos \omterm{prop:g12:adaptive-immunity:antibodies}{anticorpos}; um vírus
morto treina sobretudo o braço dos \omterm{prop:g12:adaptive-immunity:antibodies}{anticorpos}.

\textbf{4.} Os \omterm{prop:g12:adaptive-immunity:antibodies}{anticorpos} já presentes barram boa parte do vírus; os
\omterm{def:g12:adaptive-immunity:antigen}{linfócitos} B e T de memória se multiplicam em dois dias; as \omterm{def:g10:cells-common-unit:cell}{células}
infectadas são mortas antes de o vírus se espalhar. A infecção é
interrompida antes dos sintomas.

\textbf{5.} Três dias de multiplicação silenciosa, depois febre,
erupção cutânea e doença enquanto a resposta primária se constrói ao
longo de uma semana; recuperação por volta do dia dez, com \omterm{prop:g12:adaptive-immunity:selection}{células de
memória} e imunidade para a vida toda --- se ele sobreviver às complicações.

\textbf{6.} $1 - 1/15 \approx 93\%$.

\textbf{7.} $15 \times 0.10 = 1.5$: cada caso infecta mais de uma outra
pessoa, e assim a doença consegue se espalhar. Não protegida.

\textbf{8.} $10\% \times \num{50000} = 5000$ não imunes, dos quais 1500
são os bebês.

\textbf{9.} 2500 suscetíveis; $15 \times 0.05 = 0.75$: menos de um, a
doença se extingue.

\textbf{10.} Os bebês não podem ser vacinados, e assim a única proteção
deles é que os adultos em volta, sendo imunes, não deixam ao vírus
caminho nenhum para chegar até eles.

\textbf{11.} Geração 3: $1.5^3 \approx 3.4$ casos; geração 6:
$1.5^6 \approx 11$.

\textbf{12.} $1 + 1.5 + 2.25 + 3.4 + 5.1 + 7.6 + 11.4 \approx 32$ casos.

\textbf{13.} Cerca de 30\% de 32: uns 10 bebês.

\textbf{14.} Cada caso retira uma pessoa suscetível e acrescenta uma
imune, e assim a fração de suscetíveis cai e com ela o número que cada
caso infecta; com 5000 suscetíveis o surto só desaceleraria depois de
centenas de casos, ao longo de meses --- a não ser que a vacinação intervenha.

\textbf{15.} Com 95\%: $0.75^3 \approx 0.4$ e $0.75^6 \approx 0.18$
casos; total de cerca de $1 + 0.75 + 0.56 + \dots \approx 4$ casos. Dez
vezes menos casos, e o surto se extingue sozinho.

\textbf{16.} 800 em 8 anos: 100 por microlitro por ano, cerca de 8 por mês.

\textbf{17.} Os \omterm{def:g12:adaptive-immunity:antigen}{linfócitos} B precisam das citocinas das \omterm{def:g10:cells-common-unit:cell}{células}
auxiliares para virar plasmócitos, e os \omterm{def:g12:adaptive-immunity:antigen}{linfócitos} T citotóxicos
precisam delas para se multiplicar: com auxiliares de menos, nenhum dos
braços é autorizado, embora as \omterm{def:g10:cells-common-unit:cell}{células} deles existam.

\textbf{18.} Sim: as \omterm{prop:g12:adaptive-immunity:selection}{células de memória} ainda precisam das citocinas
auxiliares para se expandir numa resposta secundária; sem coordenação a
memória não consegue ser usada, e o sarampo, inofensivo para um adulto
saudável, pode ser fatal.

\textbf{19.} A \omterm{def:g12:selection-drift-speciation:population}{população} de auxiliares volta a crescer a partir das
\omterm{def:g10:cells-common-unit:cell}{células} sobreviventes e da medula, e as \omterm{prop:g12:adaptive-immunity:selection}{células de memória} dos outros
braços nunca foram destruídas --- o repertório está intacto assim que a
coordenação volta. O vírus persiste em \omterm{def:g10:cells-common-unit:cell}{células} escondidas, e assim
parar os medicamentos o deixa recomeçar.

\textbf{20.} 93\% de cobertura; cerca de 32 casos a partir de um
viajante com 90\%; o \omterm{prop:g12:adaptive-immunity:tcells}{linfócito T auxiliar}.
\end{solution}
